\chapter{Conclusions}

This project investigated randomized Nystr\"om approximation for kernel ridge regression. The main objective was to determine how the column sampling strategy and the number of sampled columns affect prediction accuracy, approximation quality and computational cost. Uniform sampling, approximate leverage score sampling and recursive ridge leverage score sampling were compared against a full KRR baseline. The final experiment used five independent column selections on the Superconductivity, CASP protein structure and CT Slice Localization datasets, with the same training size and a fixed train/test split for each dataset.

The first research question asked which Nystr\"om column sampling strategies give good performance in terms of prediction accuracy and computational cost. In the experiments, uniform sampling was the strongest practical baseline. It has essentially no sampling cost and its prediction results were very close to those of the more informed methods. Approximate leverage sampling did not consistently improve on uniform sampling. The diagnostic experiments showed that its estimated scores can be noisy and depend strongly on the quality of the initial uniform sample. Therefore the additional score estimation step did not reliably produce better landmark sets.

Recursive RLS provides a more systematic and theoretically motivated way of estimating column importance \citep{musco2017recursive}. It could produce a small kernel approximation error when sufficiently many columns were selected, but it also had the largest sampling cost and did not give a clear improvement in the final prediction-runtime trade off. In particular, the smaller kernel approximation error obtained by recursive RLS at large column counts did not imply the smallest coefficient error or test RMSE. Its sampling parameter $\lambda_{\mathrm{RLS}}$ also required careful selection because it controls the effective dimension and therefore the number of returned landmarks \citep{alaoui2015fast,musco2017recursive}. The large values needed in the experiment are related to the scale of the RBF kernel matrix. Since these values were selected independently from $\lambda_{\mathrm{KRR}}$, the recursive results should be interpreted empirically rather than as a direct verification of the theoretical guarantee.

The second research question asked how large the sampled column count $l$ needs to be and how this varies across datasets, sampling strategies and for recursive RLS, $\lambda_{\mathrm{RLS}}$. No single value of $l$ is sufficient for every required level of accuracy or every dataset. Small landmark sets gave substantial runtime reductions, but their prediction errors were noticeably larger than full KRR. Increasing $l$ generally reduced the kernel error, coefficient error and test RMSE. The effect of increasing $l$ differed significantly between datasets. CASP approached full KRR with relatively few columns, Superconductivity approached it more gradually, and CT Slice Localization remained far from full KRR over the tested range. Within each dataset, the required value of $l$ was not consistently reduced by using the more informed sampling methods. The kernel matrix and its spectral decay seemed to have a stronger effect on the required landmark count than the choice between the three tested sampling strategies. This is consistent with theoretical analyses that relate the required approximation rank to the spectrum and effective dimension of the kernel matrix \citep{bach2013sharp,rudi2015less}.

The third research question asked what the accuracy/cost trade off looks like as $l$ increases and whether a reasonable practical approach can be recommended. The main Superconductivity experiment showed a clear saturation pattern. The first increases in $l$ produced large improvements in prediction accuracy for a relatively small increase in runtime. Closing the final part of the gap to full KRR was much more expensive. In the main experiment, a point around ten percent relative excess test RMSE required approximately $4$ to $5$ seconds, or roughly one fifth of the full KRR total reference time. Reducing the gap to approximately $5\%$ required around $10$ seconds, corresponding to approximately $40\%$ of the full KRR total reference time. This provides a reasonable compromise when a small loss of prediction accuracy is acceptable. The additional datasets show that the position of this trade off is not universal as one would expect. CASP gave a substantially more favourable trade off, while CT Slice Localization was not close to the full KRR prediction region within the tested range.

We also saw that column selection was not the dominant cost at large $l$. Once the landmarks had been chosen, every sampling method used the same factorized solver with dominant cost $\mathcal{O}(nl^2+l^3)$ \citep{gittens2016revisiting}. Therefore even a theoretically better sampling method is only useful in practice if it reaches the required prediction quality with a substantially smaller column selection set. When $l$ becomes a significant fraction of $n$, the common solve dominates and the differences between the sampling costs become relatively unimportant.

Overall, the practical recommendation from this project is to begin with uniform sampling. It is simple, cheap and was competitive throughout the experiments. A more informed sampling method is worth it only when the kernel matrix has sufficiently nonuniform spread of directions in columns and the approximation method can use this structure to achieve the same prediction quality with less landmarks. This agrees with previous results showing that leverage-based sampling offers less advantage when the leverage scores are relatively uniform \citep{gittens2016revisiting,musco2017recursive}. In the experiments we did not observe this advantage consistently.

\section{Limitations and future work}

The final experiments used three datasets which gives a stronger comparison than a single dataset experiment but it is still too small to establish a universal ranking of the sampling methods. In addition, only one fixed train/test split was used for each dataset. The five repetitions therefore measure variation caused by randomized column selection, but they do not measure variation caused by changing the training and test observations. Lastly the amount of training data we used were conservative in consideration of computing power limitations. Future experiments should repeat the comparison over several independent train/test splits, a wider collection of datasets and a much larger training size.

The same RBF kernel and the same kernel rule $\gamma=\frac{1}{d}$ were used for all three datasets. This gives a controlled comparison, but it does not mean that the KRR model was optimally tuned for every dataset. In particular, the poor approximation results on CT Slice Localization may partly be the result of a kernel configuration that retained many important spectral directions. Future work should tune $\gamma$, $\lambda_{\mathrm{KRR}}$ and $l$ in combination for each dataset and then investigate relative changes.

A further limitation is the runtime measurements. Runtime depends on the experimental machine and its current workload. Background processes and other system level effects can change the measured runtime significantly even when the data and algorithms are unchanged. The reported runtimes and their corresponding analysis should therefore be interpreted as relative results rather than universal results. The asymptotic growth with $l$ and the observation that the common factorized solve becomes dominant are much more transferable than the exact runtime in seconds. A better benchmark experiment could be run on an idle and advanced machine, include warm up runs (first runs being slow was observed as outliers), repeat every measurement several times and ideally report the hardware and software configuration.

The recursive sampling parameter $\lambda_{\mathrm{RLS}}$ was selected separately from $\lambda_{\mathrm{KRR}}$ and chosen to produce computationally manageable column counts. Using a value near $\lambda_{\mathrm{KRR}}=0.01$ would retain a very large effective dimension and cause recursive RLS to select a large fraction of the training set (most of the time the entire set). Consequently the recursive experiments provide an empirical comparison of the algorithm at different returned column counts, but they do not directly verify the theoretical KRR guarantee associated with a particular coupling between regularization and column sampling. Future work should study this coupling in more detail including the effect of kernel selection and training data size.

Most importantly, the results show that improving sampling strategy alone does not resolve the runtime bottleneck. Once a large value of $l$ is required, the direct factorized solver has dominant cost $\mathcal{O}(nl^2+l^3)$ regardless of how the columns were selected. Future work could investigate iterative Nystr\"om solvers based on preconditioned conjugate gradients rather than only the direct factorized solver we used in this project \citep{rudi2017falkon}.

Finally, the experiments showed that a smaller Frobenius kernel error did not necessarily lead to a smaller coefficient error or test RMSE, suggesting that prediction-aware landmark selection may be a useful direction for future research \citep{cortes2010impact,alaoui2015fast}. Although prediction error was evaluated in this project with the two nonuniform sampling strategies using information from the kernel matrix, none of the three methods used validation prediction error when selecting landmarks. Future work could use a three-way training/validation/test split, where the validation set guides the choice of $l$ while the test set remains reserved for final evaluation. Such a rule should stop increasing $l$ when the additional validation error reduction no longer justifies the additional runtime.